\enquestion{9.42}{
    A melting point test of $n = 15$ samples of a binder used in manufacturing a rocket propellant resulted in $\overline{x} = 68^\circ\text{C}$.
    Assume that the melting point is normally distributed with $\sigma = 0.83^\circ\text{C}$.
    \begin{enumerate}[label=(\alph*)]
        \item Test $H_{0}: \mu = 68.3$ versus $H_{1}: \mu \ne 68.3$ using $\alpha = 0.01$.
        \item What is the $p$-value for this test?
        \item What is the $\beta$-error if the true mean is $\mu = 65.5$?
    \end{enumerate}
}
\enanswer{
    \begin{enumerate}[label=(\alph*)]
        \item Notice that the mean melting point of $15$ samples is normally distributed: $\overline{X} \sim N(\mu=\text{?}, \sigma=\frac{0.83}{\sqrt{15}})$.
        We perform a two-sided hypothesis test with $H_{0}: \mu = 68.3$ versus $H_{1}: \mu \ne 68.3$.
        Under the null hypothesis, we assume that $\mu = 68.3$ and check how probable is the sample mean value $\overline{x} = 68$.
        The critical region is of the form $Z = (-\infty, C_L] \cup [C_R, \infty)$, where 
        \begin{align*}
            C_L &= \invnorm(\opp=\alpha/2=0.005, \mu=68.3, \sigma=\frac{0.83}{\sqrt{15}}) \approx 67.7480. \\
            C_R &= \invnorm(\opp=1-\alpha/2=0.995, \mu=68.3, \sigma=\frac{0.83}{\sqrt{15}}) \approx 68.8520.
        \end{align*}
        The sample mean $\overline{x} = 68$ does not lie in the critical region, hence we fail to reject $H_0$.
        Based on this sample, there is not evidence to reject the claim that the true mean is $\mu=68.3$.\\[2ex]

        \textbf{Alternative:} we can also consider the $z$-score of the sample mean and compute the critical region in terms of the standard normal distribution.
        In that case, the $z$-score of the sample mean equals $z = \frac{\overline{x} - \mu_0}{\sigma/\sqrt{n}} = \frac{68-68.3}{0.83/\sqrt{15}} \approx -1.3999$.
        The critical region is of the form $Z = (-\infty, C_L] \cup [C_R, \infty)$, where 
        \begin{align*}
            C_L &= \invnorm(\opp=\alpha/2=0.005, \mu=0, \sigma=1) \approx -2.5758. \\
            C_R &= \invnorm(\opp=1-\alpha/2=0.995, \mu=0, \sigma=1) \approx 2.5758.
        \end{align*}
        The $z$-score of the sample mean $z = -1.3999$ does not lie in the critical region, hence we fail to reject $H_0$.
        Based on this sample, there is not evidence to reject the claim that the true mean is $\mu=68.3$.\\[2ex]

        \item First, note that the sample mean $\overline{x} = 68$ is smaller than the assumed null mean $\mu_0 = 68.3$.
        The $p$-value (the probability that we get an observed sample mean further from $68.3$ than $\mu_0=68$) can therefore be computed as follows:
        \begin{align*}
            p = 2 \cdot P(\overline{X} \leq \overline{x})   &= \normalcdf(\lb=-\GRinfinity, \ub=68, \mu=68.3, \sigma=\frac{0.83}{\sqrt{15}}) \\
                                                            &\approx 0.1616.
        \end{align*}
        \item The $\beta$-error (probability of not rejecting $H_0$ when $H_1$ is true) can be computed by the probability of observing a sample mean in the acceptance region,
        given the true mean $\mu = 65.5$:
        \begin{align*}
            \beta = P(H_0 \mid H_1)    &= P(C_L \leq \overline{X} \leq C_R \mid \mu=65.5) \\
                                    &= P(67.7480 \leq \overline{X} \leq 68.8520 \mid \mu=65.5) \\
                                    &= \normalcdf(\lb=67.7480, \ub=68.8520, \mu=65.5, \sigma=\frac{0.83}{\sqrt{15}}) \\
                                    &\approx 4.9268 \cdot 10^{-26}.
        \end{align*}
    \end{enumerate}
}

\enquestion{9.47}{
    Medical researchers have developed a new artificial heart constructed primarily of titanium and plastic.
    The heart will last and operate almost indefinitely once it is implanted in the patient's body, but the battery pack needs to be recharged about every four hours.
    A random sample of $50$ battery packs is selected and subjected to a life test.
    The average life of these batteries is $4.05\text{ h}$.
    Assume that battery life is normally distributed with standard deviation $\sigma = 0.2\text{ h}$.
    \begin{enumerate}[label=(\alph*)]
        \item Is there evidence to support the claim that mean battery life exceeds $4\text{ h}$?
        Use $\alpha = 0.05$.
        \item What is the $p$-value for the test in part (a)?
        \item Compute the power of the test if the true mean battery life is $4.5\text{ h}$.
    \end{enumerate}
}
\enanswer{
    \begin{enumerate}[label=(\alph*)]
        \item Notice that the mean battery life of $50$ battery packs is normally distributed: $\overline{X} \sim N(\mu=\text{?}, \sigma=\frac{0.2}{\sqrt{50}})$.
        We perform a right-sided hypothesis test with $H_{0}: \mu = 4$ versus $H_{1}: \mu > 4$.
        Under the null hypothesis, we assume that $\mu = 4$ and check how probable is the sample mean value $\overline{x} = 4.05$.
        The critical region is of the form $Z = [C_R, \infty)$, where 
        \begin{align*}
            C_R &= \invnorm(\opp=1-\alpha=0.95, \mu=4, \sigma=\frac{0.2}{\sqrt{50}}) \approx 4.0465.
        \end{align*}
        The sample mean $\overline{x} = 4.05$ lies inside the critical region, hence we reject $H_0$.
        Based on this sample, there is enough evidence to support the claim that the mean battery life indeed exceeds four hours.
        \item First, note that we consider a right-sided hypothesis test, therefore the $p$-value
        (the probability that we will observe a sample mean even larger than $4.05$) can be computed as follows:
        \begin{align*}
            p = P(\overline{X} \geq \overline{x}) = \normalcdf(\lb=4.05, \ub=\GRinfinity, \mu=4, \sigma=\frac{0.2}{\sqrt{50}}) \approx 0.0385.
        \end{align*}
        \item The power of a hypothesis test is $1 - \beta$, where $\beta$ is the type II error probability (probability of not rejecting $H_0$ when $H_1$ is true).
        With other words, the power is the probability of observing a sample mean in the acceptance region,
        given the true (alternative) mean $\mu = 4.5$:
        \begin{align*}
            1 - \beta = P(H_1 \mid H_1) &= P(\overline{X} \geq C_R \mid \mu=4.5) \\
                                     &= P(\overline{X} \geq 4.0465 \mid \mu=4.5) \\
                                     &= \normalcdf(\lb=4.0465, \ub=\GRinfinity, \mu=4.5, \sigma=\frac{0.2}{\sqrt{50}}) \\
                                     &= 1.
        \end{align*}
    \end{enumerate}
}

\enquestion{9.51}{
    For the hypothesis test $H_{0}: \mu = 7$ against $H_{1}: \mu \ne 7$ with variance unknown and $n = 25$,
    approximate the $p$-value for each of the following test statistics.
    
    \begin{center}
        \begin{tabular}{*{3}{m{0.3\textwidth}}}
            (a) $t_{0} = 2.05$ & (b) $t_{0} = -1.84$ & (c) $t_{0} = 0.4$
        \end{tabular}
    \end{center}
}
\enanswer{
    In all three cases, we consider a two-sided hypothesis test for the mean $\mu$ of a normal distribution (with unknown $\sigma$).
    Therefore, we need to resort to the $t$-distribution.

    The test statistic we are working with is 
    \[
        T = \frac{\overline{X} - \mu}{S / \sqrt{n}},
    \]
    where $\overline{X}$ is the sample mean and $S$ is the sample standard deviation.
    This test statistic follows a $t$-distribution with $\df = n - 1 = 24$ degrees of freedom.

    As we consider two-sided tests, the $p$-values of given $t_0$-values can be computed as follows:
    \begin{align*}
        p 
        = 
        \begin{cases} 
            2 \cdot P(T \geq t_{0}) = \tcdf(\lb=t_0, \ub=\GRinfinity, \df=24), & \text{ if } t_0 \geq 0 \\
            2 \cdot P(T < t_0) = \tcdf(\lb=-\GRinfinity, \ub=t_0, \df=24), & \text{ otherwise.}
        \end{cases}
    \end{align*}

    \begin{enumerate}[label=(\alph*)]
        \item Filling in $t_0 = 2.05$ gives
        \[
            p = 2 \cdot P(T \geq 2.05) = \tcdf(\lb=2.05, \ub=\GRinfinity, \df=24) \approx 0.0514.
        \]
        \item Filling in $t_0 = -1.84$ gives
        \[
            p = 2 \cdot P(T < -1.84) = \tcdf(\lb=-\GRinfinity, \ub=-1.84, \df=24) \approx 0.0782.
        \]
        \item Filling in $t_0 = 0.4$ gives
        \[
            p = 2 \cdot P(T \geq 0.4) = \tcdf(\lb=0.4, \ub=\GRinfinity, \df=24) \approx 0.6927.
        \]
    \end{enumerate}
}

\enquestion{9.59}{
    A 1992 article in the \textit{Journal of the American Medical Association} (\enquote{A Critical Appraisal of 98.6 Degrees F, the Upper Limit of the Normal Body Temperature, and Other Legacies of Carl Reinhold August Wunderlich}) reported body temperature, gender, and heart rate for a number of subjects.
    The body temperatures for $25$ female subjects follow: $97.8$, $97.2$, $97.4$, $97.6$, $97.8$, $97.9$, $98.0$, $98.0$, $98.0$, $98.1$, $98.2$, $98.3$, $98.3$, $98.4$, $98.4$, $98.4$, $98.5$, $98.6$, $98.6$, $98.7$, $98.8$, $98.8$, $98.9$, $98.9$, and $99.0$.
    \begin{enumerate}
        \item[(a)] Test the hypothesis $H_{0}: \mu = 98.6$ versus $H_{1}: \mu \ne 98.6$, using $\alpha = 0.05$.
        Find the $p$-value.
        \item[(e)] Explain how the question in part (a) could be answered by constructing a two-sided confidence interval on the mean female body temperature.
    \end{enumerate}
}
\enanswer{
    \begin{enumerate}
        \item[(a)] 
        We consider a two-sided $t$-test for the mean $\mu$ of a normal distribution, as the standard deviation $\sigma$ is unknown.
        First, we need to compute the sample mean $\overline{x}$ and the sample standard deviation $s$ of the given sample:
        \begin{align*}
            \overline{x} = \frac{x_1+x_2+\ldots+x_n}{n} = \frac{ 97.8 + 97.2 + \ldots + 99.0 }{ 25 } \approx 98.2640.
        \end{align*}
        \begin{align*}
            s^2 &= \frac{ (x_1 - \overline{x})^2 + (x_2 - \overline{x})^2 + \ldots + (x_n - \overline{x})^2 }{ n - 1 } \\
                &= \frac{ (97.8 - 98.2640)^2 + (97.2 - 98.2640)^2 + \ldots + (99.0 - 98.2640)^2 }{ 25 - 1 } \\
                &\approx 0.2324.
        \end{align*}
        \begin{align*}
            s = \sqrt{ s^2 } = \sqrt{ 0.2324 } \approx 0.4821.
        \end{align*}
    
        The test statistic for the $t$-test is then computed as
        \[
            t_0 = \frac{\overline{x} - \mu}{s / \sqrt{n}} = \frac{98.2640 - 98.6}{0.4821 / \sqrt{25}} \approx -3.4849.
        \]

        The critical region ($\alpha = 0.05$, two-sided $t$-test) is of the form $Z = (-\infty, C_L] \cup [C_R, \infty)$, where:
        \begin{align*}
            C_L &= \invt(\opp=\alpha/2 = 0.025, \df=24) \approx -2.0639, \\
            C_R &= \invt(\opp=1-\alpha/2 = 0.975, \df=24) \approx 2.0639.
        \end{align*}
        The observed test statistic $t_0$ lies inside the critical region, therefore we reject $H_0$.

        Alternatively, we could have computed the $p$-value corresponding to $t_0$:
        \begin{align*}
            p &= 2 \cdot P(T \leq -3.4849)  &= 2 \cdot \tcdf(\lb=-\GRinfinity, \ub=-3.4849, \df=24) \\
                                            &\approx 0.0019.
        \end{align*}
        As $p < \alpha$, we again decide to reject $H_0$.
        There is enough evidence to assume that the mean female body temperature significantly differs from $98.6^\circ F$.
        \item[(e)] A two-sided confidence interval for the mean $\mu$ can be computed as follows:
        \begin{align*}
            &[\overline{x} - t_{\alpha/2;n-1} \cdot \frac{s}{\sqrt{n}}, \overline{x} + t_{\alpha/2;n-1} \cdot \frac{s}{\sqrt{n}}] \\
        \end{align*}
        Here, $\overline{x} = 98.2640$, $s = 0.4821$, $n=25$ and
        \[
            t_{\alpha/2,n-1} = \invt(\opp=1-\alpha/2=0.975, \df=n-1=24) \approx 2.0639.
        \]
        Filling in this information in the formula gives an interval $[98.0650, 98.4630]$.
        Notice that $98.6$ does not lie inside this confidence interval, hence it is very unlikely that $98.6$ is the true value of the mean $\mu$.
        Therefore, we can again decide to reject $H_0$.
    \end{enumerate}
}

\enquestion{9.60}{
    Cloud seeding has been studied for many decades as a weather modification procedure (for an interesting study of this subject, see the article in \textit{Technometrics}, \enquote{A Bayesian Analysis of a Multiplicative Treatment Effect in Weather Modification,} Vol. 17, pp. 161–166).
    The rainfall in cubic-meters from $20$ clouds that were selected at random and seeded with silver nitrate follows: $22.2 \times 10^{3}$, $37.9 \times 10^{3}$, $24.4 \times 10^{3}$, $33.4 \times 10^{3}$, $27.5 \times 10^{3}$, $23.2 \times 10^{3}$, $39.2 \times 10^{3}$, $28.9 \times 10^{3}$, $26.1 \times 10^{3}$, $34.4 \times 10^{3}$, $39.3 \times 10^{3}$, $33.4 \times 10^{3}$, $30.8 \times 10^{3}$, $30.4 \times 10^{3}$, $33.2 \times 10^{3}$, $26.9 \times 10^{3}$, $36.2 \times 10^{3}$, $42.9 \times 10^{3}$, $32.9 \times 10^{3}$, and $38.9 \times 10^{3}$.
    \begin{enumerate}
        \item[(a)] Can you support a claim that the mean rainfall from seeded clouds exceeds \SI{30e3}{\cubic\metre}?
        Use $\alpha = 0.01$.
        Find the $p$-value.
        \item[(e)] Explain how the question in part (a) could be answered by constructing a one-sided confidence bound on the mean rainfall in cubic metres.
    \end{enumerate}
}
\enanswer{
    \begin{enumerate}
        \item[(a)]
        We consider a right-sided $t$-test for the mean $\mu$ of a normal distribution, as the standard deviation $\sigma$ is unknown.
        First, we need to compute the sample mean $\overline{x}$ and the sample standard deviation $s$ of the given sample:
        \begin{align*}
            \overline{x} = \frac{x_1+x_2+\ldots+x_n}{n} = \frac{ 22.2 + 37.9 + \ldots + 38.9 }{ 20 } \approx 32.1050.
        \end{align*}
        \begin{align*}
            s^2 &= \frac{ (x_1 - \overline{x})^2 + (x_2 - \overline{x})^2 + \ldots + (x_n - \overline{x})^2 }{ n - 1 } \\
                &= \frac{ (22.2 - 32.1050)^2 + (37.9 - 32.1050)^2 + \ldots + (38.9 - 32.1050)^2 }{ 20 - 1 } \\
                &\approx 34.8226.
        \end{align*}
        \begin{align*}
            s = \sqrt{ s^2 } = \sqrt{ 34.8226 } \approx 5.9011.
        \end{align*}
    
        The test statistic for the $t$-test is then computed as
        \[
            t_0 = \frac{\overline{x} - \mu}{s / \sqrt{n}} = \frac{32.1050 - 30}{5.9011 / \sqrt{20}} \approx 1.5953.
        \]

        The critical region ($\alpha = 0.01$, right-sided $t$-test) is of the form $Z = [C_R, \infty)$, where:
        \begin{align*}
            C_R &= \invt(\opp=1-\alpha = 0.99, \df=19) \approx 2.5395.
        \end{align*}
        The observed test statistic $t_0 \approx 1.5953$ does not lie inside the critical region, therefore we fail to reject $H_0$.

        Alternatively, we could have computed the $p$-value corresponding to $t_0$:
        \begin{align*}
            p &= P(T \geq 1.5953) = \normalcdf(\lb=1.5953, \ub=\GRinfinity, \df=19) \approx 0.0636.
        \end{align*}
        As $p > \alpha = 0.05$, we again fail to reject $H_0$.
        There is not enough evidence to assume that the mean rainfall from a random cloud
        significantly exceeds \SI{30e3}{\cubic\metre}.
        
        \item[(e)] An alternative way of answering part (a) is through a one-sided confidence interval for the mean $\mu$.
        As we consider a right-sided hypothesis test, this means that the confidence interval is of the form $[\overline{x} - t_{\alpha,n-1} \cdot s / \sqrt{n}, \infty)$,
        In this case, $n = 20$ and $\alpha = 0.01$, so
        \begin{align*}
            t_{\alpha,n-1} = \invt(\opp=1 - \alpha=0.99, \df=n-1=19) \approx 2.5395.
        \end{align*}
        The one-sided confidence interval is therefore equal to
        \[
            [\overline{x} - t_{\alpha,n-1} \cdot s / \sqrt{n}, \infty) = [32.1050 - 2.5395 \cdot 5.9011 / \sqrt{20}, \infty) = [28.7540,\infty).
        \]
        Notice that $\mu = 30$ is inside this confidence interval, hence we fail to reject $H_0$.
        There is not enough evidence to assume that the mean rainfall from a random cloud
        significantly exceeds \SI{30e3}{\cubic\metre}.
    \end{enumerate}
}

\enquestion{9.126}{
    Consider the following computer output.

    \begin{center}
        \begin{tabular}{*{7}{C{0.11428571428\textwidth}}}
            \toprule
                \multicolumn{7}{l}{\textbf{One-Sample $Z$:}} \\                
                \multicolumn{7}{l}{Test of $\mu = 26$ vs $\mu>26$} \\
                \multicolumn{7}{l}{The assumed standard deviation $= 1.7$} \\   
                Variable & $n$ & Mean & StDev & SE Mean & $Z$ & $p$ \\
                $X$ & ? & $26.541$ & $2.032$ & $0.401$ & ? & ? \\                  
            \bottomrule
        \end{tabular}
    \end{center}

    \begin{enumerate}[label=(\alph*)]
        \item Fill in the missing information.
        \item Is this a one-sided or a two-sided test?
        \item What are your conclusions if $\alpha = 0.05$?
        \item Find a \SI{95}{\percent} two-sided CI on the mean.
    \end{enumerate}
}
\enanswer{
    \begin{enumerate}[label=(\alph*)]
        \item From the given information, we can derive the following missing descriptive statistics:
        \begin{itemize}
            \item Notice that the standard error for the mean (SE Mean) follows from the formula
            \[
                \text{SE Mean} = \frac{\sigma}{\sqrt{n}},
            \]
            or alternatively $n = \left(\frac{\sigma}{\text{SE Mean}}\right)^2 = \left(\frac{1.7}{0.401}\right) = 18$.
            \item The $z$-score corresponding to a sample with sample mean $26.541$, $\sigma = 1.7$ and $n=18$ is equal to:
            \[
                z = \frac{\overline{x} - \mu}{\sigma / \sqrt{n}} = \frac{26.541 - 26}{1.7 / \sqrt{18}} \approx 1.3502.
            \]
            \item The $p$-value corresponding to this sample is equal to
            \begin{align*}
                p &= P(\overline{X} \geq \overline{x}=26.541) \\
                &= \normalcdf(\lb=26.541, \ub=\GRinfinity, \mu=26, \sigma=1.7/\sqrt{18}) \\
                &\approx 0.0885.
            \end{align*}
        \end{itemize}
        \item This is a one-sided hypothesis test (specifically, a right-sided test), as the alternative hypothesis is of the form \enquote{$>$}.
        \item If the significance level $\alpha = 0.05$, we have that $p = 0.0885 > \alpha = 0.05$, hence we fail to reject $H_0$.
        There is not enough evidence to reject the claim that the true mean $\mu = 26$.
        \item A two-sided \SI{95}{\percent} ($\alpha = 0.05$) confidence interval for the mean $\mu$ can be computed as follows:
        \begin{align*}
            &[\overline{x} - z_{\alpha/2} \cdot \frac{\sigma}{\sqrt{n}}, \overline{x} + z_{\alpha/2} \cdot \frac{\sigma}{\sqrt{n}}]
        \end{align*}
        Here, $\overline{x} = 26.541$, $\sigma = 1.7$, $n=18$ and
        \[
            z_{\alpha/2} = \invnorm(\opp=1-\alpha/2=0.975, \mu=0, \sigma=1) \approx 1.9600.
        \]
        Filling in this information in the formula gives an interval 
        \begin{align*}
            &[\overline{x} - z_{\alpha/2} \cdot \frac{\sigma}{\sqrt{n}}, \overline{x} + z_{\alpha/2} \cdot \frac{\sigma}{\sqrt{n}}] \\
            &=[26.541 - 1.9600 \cdot \frac{1.7}{\sqrt{18}}, 26.541 + 1.9600 \cdot \frac{1.7}{\sqrt{18}}] \\
            &=[25.7556, 27.3264].
        \end{align*}
    \end{enumerate}
}

\enquestion{9.127}{
    Consider the following computer output.

    \begin{center}
        \begin{tabular}{*{8}{C{0.1\textwidth}}}
            \toprule
                \multicolumn{8}{l}{\textbf{One-Sample $T$:}} \\                
                \multicolumn{8}{l}{Test of $\mu = 100$ vs $\mu\neq 100$} \\
                Variable & $n$ & Mean & StDev & SE Mean & \SI{95}{\percent} CI & $t$ & $p$ \\
                $X$ & $16$ & $98.33$ & $4.61$ & ? & $($?,?$)$ & ? & ? \\                  
            \bottomrule
        \end{tabular}
    \end{center}
    
    \begin{enumerate}[label=(\alph*)]
        \item How many degrees of freedom are there on the $t$-statistic?
        \item Fill in the missing information.
        You may use bounds on the $p$-value.
        \item What are your conclusions if $\alpha = 0.01$?
        \item What are your conclusions if the hypothesis is $H_{0}: \mu = 100$ versus $H_{0}: \mu > 100$?
    \end{enumerate}
}
\enanswer{
    \begin{enumerate}[label=(\alph*)]
        \item The number of degrees of freedom equals $\df = n - 1$, where $n = 16$ is the sample size.
        Therefore, $\df = 15$.
        \item From the given information, we can derive the following missing descriptive statistics:
        \begin{itemize}
            \item Notice that the standard error for the mean (SE Mean) follows from the formula
            \[
                \text{SE Mean} = \frac{\text{StDev}}{\sqrt{n}} = \frac{4.61}{\sqrt{16}} \approx 1.1525.
            \]
            \item A two-sided \SI{95}{\percent} ($\alpha = 0.05$) confidence interval for the mean $\mu$ can be computed as follows:
            \begin{align*}
                &[\overline{x} - t_{\alpha/2,n-1} \cdot \frac{\sigma}{\sqrt{n}}, \overline{x} + t_{\alpha/2,n-1} \cdot \frac{\sigma}{\sqrt{n}}]
            \end{align*}
            Here, $\overline{x} = 98.33$, $s = 4.61$, $n=16$ and
            \[
                t_{\alpha/2,n-1} = \invt(\opp=1-\alpha/2=0.975, \df=15) \approx 2.1314.
            \]
            Filling in this information in the formula gives an interval 
            \begin{align*}
                &[\overline{x} - t_{\alpha/2,n-1} \cdot \frac{s}{\sqrt{n}}, \overline{x} + t_{\alpha/2,n-1} \cdot \frac{s}{\sqrt{n}}] \\
                &=[98.33 - 2.1314 \cdot \frac{4.61}{\sqrt{16}}, 98.33 + 2.1314 \cdot \frac{4.61}{\sqrt{16}}] \\
                &=[95.8735, 100.7865].
            \end{align*}
            \item The $t$-score corresponding to a sample with $\overline{x}=98.33$, $s = 4.61$, $\mu=100$ and $n=16$ is equal to:
            \[
                t = \frac{\overline{x} - \mu}{s / \sqrt{n}} = \frac{98.33 - 100}{4.61 / \sqrt{16}} \approx -1.4490.
            \]
            \item The $p$-value corresponding to this sample is equal to
            \[
                p = 2 \cdot P(T \leq t) = 2 \cdot \tcdf(\lb=-\GRinfinity, \ub=-1.4490, \df=15) \approx 0.1679.
            \]
        \end{itemize}
        \item If the significance level $\alpha = 0.01$, we have that $p = 0.1679 > \alpha = 0.01$, hence we fail to reject $H_0$.
        There is not enough evidence to reject the claim that the true mean $\mu = 100$.
        \item Suppose now we consider a right-sided hypothesis test.
        In that case, the critical region ($\alpha = 0.01$, right-sided $t$-test) is of the form $Z = [C_R, \infty)$, where:
        \begin{align*}
            C_R &= \invt(\opp=1-\alpha = 0.99, \df=15) \approx 2.6025.
        \end{align*}
        The observed test statistic $t_0 \approx -1.4490$ does not lie inside the critical region, therefore we fail to reject $H_0$.

        Alternatively, we could have computed the $p$-value corresponding to $t_0$:
        \begin{align*}
            p &= P(T \geq -1.4490) = \tcdf(\lb=-1.4490, \ub=\GRinfinity, \df=15) \approx 0.9160.
        \end{align*}
        As $p > \alpha = 0.01$, we again fail to reject $H_0$.
        There is not enough evidence to assume that the true mean is significantly larger than $100$.
    \end{enumerate}
}